\begin{textsample}{POS Dim 3 – persona\_grok – Score 18.00 – t118\_grok.txt}  \label{ex:f3_pos_035}
Europe 's \textbf{transport} \textbf{system}, heavily reliant on fossil fuels, is fueling a cascade of crises that affect us all.
From accelerating climate change to driving inflation and skyrocketing \textbf{energy} \textbf{prices}, it 's also propping up conflicts funded by oil exports, such as Russia 's invasion efforts.
For \textbf{households} across the continent, the \textbf{costs} of this dependency are staggering, eating into \textbf{budgets} and livelihoods.
Transport alone accounts for about 30 % of the EU 's \textbf{greenhouse} \textbf{gas} \textbf{emissions}, making it a major contributor to the climate emergency.
Beyond \textbf{emissions}, the extraction of oil causes widespread environmental destruction, while so-called \textbf{solutions} like agrofuels often prove to be false promises that fail to address the root \textbf{problems}.
It 's \textbf{time} to drain the oil from our mobility \textbf{systems}.
By \textbf{shifting} towards affordable \textbf{public} \textbf{transport} and active \textbf{options} like walking and \textbf{cycling}, we can build a cleaner, more equitable future.
The goal is clear : fully decarbonize \textbf{transport} by 2040, steering clear of inefficient \textbf{alternatives} that only delay real progress.
To make this transition fair, we must \textbf{end} subsidies for polluters and provide strong support for workers affected by the \textbf{shift}, ensuring no one is left behind in the move to sustainable mobility.

% matched lemmas: alternative, budget, cost, cycle, emission, end, energy, gas, greenhouse, household, option, price, problem, public, shift, solution, system, time, transport
\end{textsample}
